\chapter{Scaling phase diagrams for charge, particles, and local dimension}
\label{ch:scaling-phases}
\index{particle-number scaling}
\index{quasipolynomial simulation}
\index{subexponential simulation}

\section{Exact sector size is a complexity phase variable}
The statement ``fixed particle number is polynomial'' has a sharp extension. Let the particle number depend on the system size:
\[
k=k(n).
\]
The explicit-amplitude simulation cost of one fermionic or qubit weight packet is governed by
\[
D_{n,k}=\binom nk.
\]
For $1\le k\le n$, the standard bounds
\begin{equation}
\left(\frac nk\right)^k
\le
\binom nk
\le
\left(\frac{en}{k}\right)^k
\label{eq:binomial-basic-bounds}
\end{equation}
show that the growth changes continuously between polynomial and exponential regimes.

Because
\[
\binom nk=\binom n{n-k},
\]
write
\[
r(n)=\min\{k(n),n-k(n)\}.
\]

\begin{theorem}[Explicit fixed-charge phase diagram]
\label{thm:charge-phase-diagram}
For explicit amplitude storage in a weight-$k(n)$ packet:
\begin{enumerate}
\item the packet dimension is polynomial in $n$ if $r(n)=O(1)$;
\item the packet dimension is superpolynomial whenever $r(n)\to\infty$;
\item if $r(n)=O(\log n)$, then
\[
D_{n,k}=2^{O((\log n)^2)},
\]
so the explicit packet simulation is quasipolynomial;
\item if $r(n)=O(n^\alpha)$ for a fixed $0<\alpha<1$, then
\[
D_{n,k}=2^{O(n^\alpha\log n)},
\]
which is subexponential;
\item if $k(n)/n\to p\in(0,1)$, then
\[
\log_2 D_{n,k}
=
nH_2(k(n)/n)+O(\log n)
=
nH_2(p)+o(n),
\]
where $H_2$ is binary entropy, and the packet remains exponential.
\end{enumerate}
\end{theorem}

\begin{proof}
The upper bounds follow from \cref{eq:binomial-basic-bounds}. If $r=O(1)$, then $\binom nr=O(n^r)$. Conversely, suppose $r\to\infty$ and $r\le n/2$. If $r\le\sqrt n$, then
\[
\log\binom nr
\ge
r\log(n/r)
\ge
\frac r2\log n,
\]
whose ratio to $\log n$ diverges. If $r>\sqrt n$, then
\[
\log\binom nr
\ge
r\log2
\ge
\sqrt n\log2,
\]
which also dominates every constant multiple of $\log n$. For the final assertion, $p_n=k(n)/n$ stays in a compact subinterval of $(0,1)$. Stirling's formula gives
\[
\log_2\binom n{k(n)}
=nH_2(p_n)-\frac12\log_2\!\bigl(2\pi n p_n(1-p_n)\bigr)+O(1/n).
\]
The logarithmic term is $O(\log n)$, uniformly on that subinterval. Continuity of $H_2$ gives $nH_2(p_n)=nH_2(p)+o(n)$. Replacing $p_n$ by its limit inside an $O(\log n)$ estimate would require a convergence rate, which is not assumed here.
\end{proof}

\begin{derivedbox}
For the explicit occupation-amplitude representation, ``few particles'' means a bounded number of particles or holes. A slowly growing charge can still improve on full state-vector scaling, but it no longer guarantees membership in $P$.
\end{derivedbox}

\section{A bounded-excitation hierarchy}
The union of the first $K$ weight sectors has dimension
\[
D_{\le K}(n)
=
\sum_{k=0}^{K}\binom nk.
\]
For fixed $K$,
\[
D_{\le K}(n)=O(n^K).
\]
If $K=O(\log n)$, the same union is quasipolynomial. This refines \cref{cor:bounded-excitation-collapse}: bounded symmetry breaking preserves polynomial simulation, while logarithmically growing excitation gives a controlled quasipolynomial regime.

\section{Bosonic and fermionic analogues}
For $N$ fermions in $m$ modes,
\[
D_{\mathrm F}(m,N)=\binom mN.
\]
For $N$ bosons in $m$ modes,
\[
D_{\mathrm B}(m,N)=\binom{m+N-1}{N}.
\]
The same binomial analysis gives:

\begin{corollary}[Few-particle scaling]
\label{cor:few-particle-scaling}
With the number of modes as the input-size parameter:
\begin{enumerate}
\item fixed $N$ gives polynomial explicit-amplitude simulation;
\item $N=O(\log m)$ gives quasipolynomial sector size;
\item $N=O(m^\alpha)$, $0<\alpha<1$, gives subexponential sector size;
\item finite-density fermionic sectors remain exponential.
\end{enumerate}
For bosons, the statements hold in the dilute regime with the same occupation-basis access assumptions.
\end{corollary}

The adjective ``polynomial'' here refers to the sector dimension and efficient gate action. It does not assert that every interacting Hamiltonian has efficiently computable matrix elements or that every compact description remains closed.

\section{Symmetric qudits with growing local dimension}
The symmetric $d$-level subspace has dimension
\[
S(n,d)
=
\binom{n+d-1}{d-1}.
\]
For fixed $d$, this is polynomial in $n$. If $d=d(n)$ grows, the same phase transition appears.

\begin{proposition}[Growing-local-dimension boundary]
\label{prop:growing-d-boundary}
In the explicit occupation basis:
\begin{enumerate}
\item fixed $d$ gives $S(n,d)=O(n^{d-1})$;
\item $d=O(\log n)$ gives $S(n,d)=2^{O((\log n)^2)}$;
\item unbounded $d$ is not covered by the fixed-$d$ polynomial-collapse theorem without an additional internal compression.
\end{enumerate}
\end{proposition}

Thus permutation symmetry removes the exponential tensor-product multiplicity only when the number of local levels is controlled, or when the occupation record itself admits a second compression.

\begin{figure}[H]
\centering
\includegraphics[width=.94\textwidth]{particle_scaling_phases.pdf}
\caption{The active dimension passes through polynomial, quasipolynomial, subexponential, and exponential regimes as particle number grows.}
\end{figure}

\section{Complexity class versus practical runtime}
A polynomial sector can still have a large exponent. Fixed $k=20$ gives $O(n^{20})$ explicit amplitudes, which is in $P$ but may be impractical. Conversely, a subexponential sector may be useful at moderate size without changing the formal complexity class.

A complete performance claim should therefore report both:
\[
\text{asymptotic class}
\qquad\text{and}\qquad
\text{actual representation size and update cost}.
\]
This is why the benchmark chapters keep wall-clock evidence separate from the dimension theorems.
